\section{Transport a resonance around the multiple zero}
\label{sec:ep-holonomy}
% BEGIN CURATED SUBJECT INDEX
\index{exceptional point}
\index{holonomy}
% END CURATED SUBJECT INDEX

\subsection{Biorthogonal Berry connection}
\label{subsec:biorthogonal-berry}
% BEGIN CURATED SUBJECT INDEX
\index{operator!adjoint and self-adjoint}
\index{biorthogonality}
\index{exceptional point}
\index{geometric phase}
\index{holonomy}
% END CURATED SUBJECT INDEX

Away from degeneracies, choose a right--left pair normalized by
\begin{equation}
  \braket{\Psi_n^{\mathrm L}(\bm\lambda)|
  \Psi_n^{\mathrm R}(\bm\lambda)}=1 ,
  \label{eq:ep-normalization}
\end{equation}
where $\bm\lambda$ denotes control parameters.  The biorthogonal Berry
connection is
\begin{equation}
  \mathcal A_n(\bm\lambda)
  =\rmi\bra{\Psi_n^{\mathrm L}(\bm\lambda)}
  \bm\nabla_{\lambda}
  \ket{\Psi_n^{\mathrm R}(\bm\lambda)} .
  \label{eq:biorthogonal-connection}
\end{equation}
For a closed loop $C$ that returns the eigenspace to itself, the geometric
holonomy is
\begin{equation}
  \gamma_n(C)=\oint_C\mathcal A_n\cdot\rmd\bm\lambda .
  \label{eq:berry-holonomy}
\end{equation}
In a non-self-adjoint problem $\gamma_n$ may be complex; its real part is a
phase and its imaginary part changes the transported amplitude.  Under a
biorthogonal gauge transformation
\begin{align}
  \ket{\Psi_n^{\mathrm R}}&\longmapsto
  \rme^{f_n}\ket{\Psi_n^{\mathrm R}},
  \notag\\
  \bra{\Psi_n^{\mathrm L}}&\longmapsto
  \bra{\Psi_n^{\mathrm L}}\rme^{-f_n},
  \label{eq:biorthogonal-gauge}
\end{align}
the connection changes by $\rmi\bm\nabla_{\lambda}f_n$, while a properly
closed holonomy is gauge invariant modulo the allowed branch and normalization
conventions.

At a second-order EP, one circuit exchanges the two eigenlines.  A scalar
Berry phase for a single eigenline is therefore defined only after enough
circuits to close the line, or by using a non-Abelian connection on the
two-dimensional root bundle.  In the standard finite-dimensional case, two
circuits return the eigenvector with a minus sign under common conventions.

\subsection{Momentum-bin regularization of the continuum}
\label{subsec:momentum-bin}
% BEGIN CURATED SUBJECT INDEX
\index{exceptional point}
\index{holonomy}
% END CURATED SUBJECT INDEX

Let continuum states satisfy
\begin{equation}
  \braket{k|k'}=\delta(k-k') .
  \label{eq:continuum-normalization}
\end{equation}
Partition a complex-scaled momentum contour into bins
$[k_n,k_{n+1}]$ with width $\Delta k_n=k_{n+1}-k_n$.  Define
\begin{equation}
  \ket{n}
  =\frac{1}{\sqrt{\Delta k_n}}
  \int_{k_n}^{k_{n+1}}\rmd k\,\ket{k} .
  \label{eq:momentum-bin-state}
\end{equation}
For nonoverlapping bins and a compatible contour pairing,
$\braket{m|n}=\delta_{mn}$.  The factor
$1/\sqrt{\Delta k_n}$ converts Dirac-delta normalization into Kronecker-delta
normalization.  It is also the source of an additional branch when the bin
width itself depends on the EP parameter.

Near a resonance--continuum coalescence, suppose the continued momentum has
the local form
\begin{equation}
  k_n(\lambda)
  =\vsub{k}{EP}
  +a_n(\lambda-\vsub{\lambda}{EP})^{1/2}
  +\cdots .
  \label{eq:k-ep-puiseux}
\end{equation}
Then
\begin{equation}
  \Delta k_n
  =(a_{n+1}-a_n)
  (\lambda-\vsub{\lambda}{EP})^{1/2}
  +\cdots .
  \label{eq:bin-width-branch}
\end{equation}
The integral in Eq.~\eqref{eq:momentum-bin-state} contributes one power of
$\Delta k_n$, while the normalization contributes its inverse square root.
Consequently the leading binned wave function scales as
\begin{equation}
  \ket{n(\lambda)}
  \propto(\lambda-\vsub{\lambda}{EP})^{1/4} .
  \label{eq:quarter-root-state}
\end{equation}
A $2\pi$ circuit in $\arg(\lambda-\vsub{\lambda}{EP})$ multiplies this
factor by $\rme^{\rmi\pi/2}$; two circuits give a minus sign, and four circuits
restore the original branch.  The quarter root combines the square-root
energy sheet with the square root required by continuum normalization
\cite{Morikawa:2025inx}.

The bin result is not an extra physical state.  It is a regulated
representation of a distributional continuum vector.  A wave-packet
regularization gives the same logic with a smooth width $\epsilon$ in place of
$\Delta k_n$.  Any claim of universality should be phrased in terms of the
closed physical holonomy after the regulator is removed, not in terms of an
individual normalization factor.

\subsection{Chern data and branch exchange}
\label{subsec:chern-branch}
% BEGIN CURATED SUBJECT INDEX
\index{topology}
\index{biorthogonality}
\index{exceptional point}
\index{Chern number}
% END CURATED SUBJECT INDEX

On a parameter surface that avoids the EP, define the curvature
\begin{equation}
  \mathcal F_n=\rmd\mathcal A_n .
  \label{eq:berry-curvature}
\end{equation}
If a closed two-cycle is well defined and the eigenline returns to itself, the
first Chern number is
\begin{equation}
  c_1=\frac{1}{2\pi}
  \int\mathcal F_n .
  \label{eq:first-chern}
\end{equation}
Near an EP, a single eigenline is branched and may not define a line bundle on
the punctured base without passing to a covering space.  One may instead work
on that cover or use the rank-two bundle of the exchanging pair.  This avoids
assigning an integer Chern number to a path that does not close in the original
bundle.  Non-Hermitian topological classifications use related biorthogonal
constructions, with additional boundary sensitivity in spatially extended
systems
\cite{Gong2018Topological,Yao2018EdgeStates,Alase:2016tqv}.

\subsection{Renormalizing a distributional holonomy}
\label{subsec:holonomy-renormalization}
% BEGIN CURATED SUBJECT INDEX
\index{monodromy}
\index{biorthogonality}
\index{exceptional point}
\index{holonomy}
\index{renormalization group}
% END CURATED SUBJECT INDEX

The continuum limit of Eq.~\eqref{eq:biorthogonal-connection} can contain
objects such as $\delta(0)$ or derivatives of a delta function if unsmeared
continuum eigenvectors are inserted directly.  These symbols are not numbers.
They stand for normalization-volume or wave-packet-width dependence.  Let
$R>0$ denote a small radius excised around the EP in parameter space, and let
$\epsilon>0$ denote a continuum smearing scale.  Write a bare holonomy as
\begin{equation}
  \vsup{\gamma}{bare}(R,\epsilon)
  =Z_{\gamma}(R,\epsilon)
  \vsup{\gamma}{phys} .
  \label{eq:holonomy-Z}
\end{equation}
The physical holonomy is defined by a renormalization condition, for example
by matching the finite-bin branch exchange.  Regulator independence requires
\begin{equation}
  \left(
  R\frac{\partial}{\partial R}
  +\epsilon\frac{\partial}{\partial\epsilon}
  \right)
  \vsup{\gamma}{phys}=0 .
  \label{eq:holonomy-rg-invariance}
\end{equation}
Equivalently, the running of $Z_\gamma$ cancels the regulator dependence of the
bare connection.  This is the EP analogue of fixing the origin monodromy in
Eq.~\eqref{eq:monodromy-Z}.

The construction is best called \textit{monodromy renormalization} or an
\textit{RG-type regulator-independence condition}.  It does not by itself
establish a Wilsonian renormalization group, a universal beta function, or a
new thermodynamic scaling law.  Its testable claim is narrower: different
continuum regulators that preserve the analytic pole and the same matching
condition should give the same closed holonomy.  Demonstrating this regulator
equivalence beyond solvable one-dimensional models is an important open
problem.

\subsection{What survives the continuum limit}
\label{subsec:ep-universal-data}
% BEGIN CURATED SUBJECT INDEX
\index{topology}
\index{exceptional point}
\index{holonomy}
% END CURATED SUBJECT INDEX

Four pieces of data have different degrees of robustness.  The location of
the continued branch point is invariant under admissible spectral
representations.  The Puiseux order is invariant unless a threshold changes
the analytic class.  The exchange permutation of spectral sheets is
topological.  The phase assigned to an individual eigenvector depends on
normalization and gauge, but a closed, matched holonomy is physical.  This
hierarchy prevents regulator-specific factors from being mistaken for new
topological invariants.

In applications, one should report the analytic family being varied, the
definition of the continuum regulator, the number of circuits required to
close the state, and the left--right normalization.  Without these data, two
  correct calculations can appear to disagree merely because they compare
  different covers or gauges.
